\documentclass[10pt]{article}

\usepackage[utf8]{inputenc}

\usepackage[T1]{fontenc}

\usepackage{amsmath}

\usepackage{amsfonts}

\usepackage{amssymb}

\usepackage[version=4]{mhchem}

\usepackage{stmaryrd}

\usepackage{bbold}

\usepackage{graphicx}

\usepackage[export]{adjustbox}

\graphicspath{ {./images/} }



\begin{document}

вонич. изоморфизм между группой Брауәра поля $K$



\[

\operatorname{Br}(K) \simeq H^{2}\left(G(\bar{K} / K), \bar{K}^{*}\right)

\]



и $\mathbb{Q} / \mathbb{Z}$



В глобальной П. к. т. роль мультипликативной группы поля играет группа классов иделей. Пусть $L / K$ - конечное расширение Галуа глобальных полей, и $I_{L}$ - группа иделей поля $L$. Группа $L^{*}$ вкладывается в $I_{L}$ в качестве дискретной подгруппы (она наз. г р у п п о й $г$ л а в в ы х и д е л е й), а факторгруппа $C_{L}=I_{L} / L^{*}$, наделенная фактортопологией, наз. г р у п п ой к л а с с о в и д е ле й. Доказывается, что $H^{1}\left(G(L / K), C_{L}\right)=1$ и $H^{2}\left(G(L / K), C_{L}\right) \simeq \mathbb{Z} / n \mathbb{Z}$, где $n=[L: K]$. Существует канонич. вложение inv: $H^{2}\left(G(L / K), \quad C_{L}\right) \rightarrow \mathbb{Q} / \mathbb{Z}$. Как и в локальной П. к. т., для любого целого $n$ определен изоморфизм (основной изоморфизм глобальной П. к. т.)



\[

\psi_{n}: H^{n-2}(G(L / K), \mathbb{Z}) \simeq H^{n}\left(G(L / K), C_{L}\right) .

\]



Для абелева расширения $L / K$ изоморфизм $\psi_{0}$ сводится к изоморфизму $\psi: G(L / K) \simeq C_{K} / N_{L / K}\left(C_{L}\right)$. Норменная подгруппа $N_{L / K}\left(C_{L}\right)$ однозначно определяет поле $L$, и наоборот, любая открытая подгруппа конечного индекса в $C_{K}$ является норменной подгруппой для нек-рого конечного абелева расширения $L$ (г л обальная те о рема с уше с т во в ания). Соотношения, аналогичные (1) и (2), остаются справедливыми и для глобальных полей. Если $K^{a b}$ - максимальное абелево расширение поля $K$, то в функциональном случае группа $G\left(K^{a b / K}\right)$ изоморфна проконечному пополнению группы $C_{K}$, а в числовом случае группа $G\left(K^{a b / K)}\right.$ изоморфна факторгруппе группы $C_{K}$ по связной компоненте.



Изоморфизмы $\varphi_{n}$ и $\psi_{n}$ согласованы. Если $L / K$ конечное расширение Галуа глобальных полей, $L_{v}-$ пополнение поля $L$ относительно нек-рой точки $v$ и $K_{v}-$ пополнение поля $K$ относительно ограничения $v$ на $K$, то существует коммутативная диаграмма



\[

\begin{gathered}

\psi_{n}: H^{n-2}(G(L / K), \mathbb{Z}) \simeq H^{n}\left(G(L / K), C_{L}\right), \\

\uparrow \text { cores } \\

\varphi_{n}: H^{n-2}\left(G\left(L_{v} / K_{v}\right), \mathbb{Z}\right) \simeq H^{n}\left(G\left(L_{v} / K_{v}\right), L^{*}\right),

\end{gathered}

\]



где отображение $f$ индуцировано вложением $L^{*}-I_{L} \rightarrow$ $\rightarrow C_{L}$ и коограничением cores. Для $n=0$ (3) дает коммутативную диаграмму



\[

\begin{aligned}

\psi: G(L / K) /[G(L / K), G(L / K)] & \simeq C_{K_{\uparrow} /} / N_{L / K}\left(C_{L}\right) \\

\varphi: G\left(L_{v} / K_{v}\right) /\left[G\left(L_{v} / K_{v}\right), G\left(L_{v} / K_{v}\right)\right] & \simeq K_{v}^{*} / N_{L_{v} / K_{v}}\left(L_{v}^{*}\right) .

\end{aligned}

\]



Диаграмма (4) позволяет получить закон разложения простых дивизоров поля $K$ в абелевом расширении $L / K$. Именно, дивизор с поля $K$ не разветвлен в $L$ (вполне распадается в $L$ ) тогда и только тогда, когда $U\left(K_{\mathfrak{c}}\right) \subset N_{L / K}\left(G_{L}\right) \quad\left(\right.$ соответственно $\left.K^{*} \subset N_{L / K}\left(G_{L}\right)\right)$.



Если $\mathfrak{c}$ - нек-рый простой дивизор поля $K$, не разветвленный в $L, v$ - точка поля $K$, соответствующая с, и $\pi$ - простой элемент поля $K_{v}$, то определен символ Артина $\left(\frac{L / K}{c}\right)=\psi^{-1}(\pi) \in G(L / K), \quad$ зависящий только от с. Элемент $\left(\frac{L / K}{c}\right)$ - это автоморфизм Фробениуса в подгруппе разложения точки $v$. Согласно т е о р еме плотности Чеботаре в а любой элемент группы $G(L / K)$ имеет вид $\left(\frac{L / K}{c}\right)$ для бесконечного числа простых дивизоров с поля $K$.



Напр., максимальное абелево неразветвленное расширение числового поля $K$ (называемое г и л ь б е рто в ы м по ле м к л а с с о в) - это поле, нормен- ная подгруппа к-рого совпадает с образом относительно проекции $I_{K} \rightarrow C_{K^{\prime}}$ группы $K^{*} \times \prod_{v} U\left(K_{v}\right)$, где $v$ пробегает все точки поля $K$. Группа $I_{K} / K^{*} \times \prod_{v} U\left(K_{v}\right)$ канонически изоморфна группе классов дивизоров $\mathrm{Cl}_{K}$ поля $K$, что дает важный изоморфизм $G(F / K) \simeq \mathrm{Cl}_{K}$. В частности, над $K$ нет неразветвленных абелевых расширений тогда и только тогда, когда поле $K$ одноклассно.



Тип разложения простого дивизора с поля $K$ в $F$ полностью определяется классом $\mathrm{c} \mathrm{s} \mathrm{Cl}_{\mathrm{c}}$. Вполне распадаются в $F$ все главные дивизоры и только оніг. Все дивизоры поля $K$ становятся главными в $F$.



Подобно тому, как П. к. т. для абелевых неразветвленных расширений можно излагать на языке группы классов дивизоров и ее подгрупп, можно дать характеризацию любого конечного абелева расширения поля $K$ в терминах группы лучевых классов по нек-рому модулю (см. Алгебрачческая пеория чисел). Существуют также обобщения П. к. т. на случай бесконечных расширений Галуа [4].



Хотя П. к. т. возникла как теория абелевых расширений, ее результаты дают важную информацию и для неабелевых расширений Галуа. Напр., на теории полей классов основано доказательство существования бесконечных башен полей классов (см. Башня полей).



Лит.: [1] Алгебраическая теория чисел, пер. с англ., М., 1969 ; [2] В е й л ь А., Основы теории чисел, пер. с англ., М., 1972; [3] К о х Х., Теория Галуа $p$-расширений, пер. с нем., М., $1973 ;$ [4] К у з в м и н Л. В., «Изв. А Н СССР. Сер. матем.»,

1969 , т. 33, в. 6, с. $1220-54$.



ПОЛИАНАЛИТИЧЕСКАЯ ФУНКЦИЯ п о р я дк $\mathbf{a} m-$ комплексная функция $w=u+i v$ действительных переменных $x$ и $y$ или, что эквивалентно, независимых комплексных переменных $z=x+i y$ и $\bar{z}=x-i y$ в плоской области $D$, представимая в виде



\[

w=f(z, \bar{z})=\sum_{k=0}^{m-1} \bar{z}^{k} f_{k}(z),

\]



где $f_{k}(z), k=0, \ldots, m-1,-$ комплексные аналитич. функции в $D$. Иначе, П. ф. $w$ порядка $m$ можно определить как функцию, имеющую в $D$ непрерывные частные производные по $x$ и $y$ или по $z$ и $\bar{z}$ до порядка $m$ включительно и удовлетворяющую всюду в $D$ обобщенному условию Коши - Римана:



\[

\frac{\partial^{m_{w}}}{\partial \bar{z}{ }^{m}}=0 .

\]



При $m=1$ получаются аналитич. Функции.



Для того чтобы функция $u=u(x, y)$ была действительной (или мнимой) частью нек-рой П. $\phi . w=u+i v$ в области $D$, необходимо и достаточно, чтобы $u$ была noлигармонической бункцией в $D$. На П. ф. переносятся с соответствующими изменениями нек-рые классич. свойства аналитич. функций (см. [1]).



П. ф. мультипорядка $m=\left(m_{1}, \ldots, m_{n}\right)$ от комплексных переменных $z_{1}, \ldots, z_{n}$ и $\bar{z}_{1}, \ldots, \bar{z}_{n}$ в области $D$ комплексного пространства $\mathbb{C}^{n}, n \geqslant 1$, наз. фуукция вида



$w==\sum_{k_{1}, \ldots, k_{n}=0}^{m_{1}-1, \ldots, m_{n}-1} \bar{z}_{1}^{k_{1}} \ldots \bar{z}_{n}^{k_{n}} \bar{f}_{k_{1} \ldots k_{n}}\left(z_{1}, \ldots, z_{n}\right)$,



где $f_{k_{1}} \ldots k_{n}-$ аналитич. функции переменных $z_{1}, \ldots$,



\begin{center}

\includegraphics[max width=\textwidth]{2023_07_08_88ce3e7f553102d6cda9g-1}

\end{center}



ПОЛИВЕКТОР, $p$-в е к т о р, векторного пространства $V$ - элемент $\not$-й внешней степени $\Lambda^{p} V$ пространства $V$ над полем $k$ (см. Внешнял алгебра). р-вектор может пониматься как кососимметризованный $p$ раз контравариантный тензор на $V$. Любая линейно независимая система векторов $x_{1}, \ldots, x_{p}$ из $V$ определяет





\end{document}